\begin{table}[H]
\centering
\captionsetup{justification=centering}
\caption{Medidas de desempenho dos modelos univariado e multivariado - total de desocupados}
\label{tab:diffviciodesoc}
{%
\renewcommand{\arraystretch}{0.8}
\scalebox{0.95}{%
\begin{tabular}{lcccc}
\toprule
& \multicolumn{2}{c}{\textbf{\makecell{Diferença relativa média \\ do erro padrão (\%)}}} & \multicolumn{2}{c}{\textbf{Vício relativo (\%)}} \\
\cmidrule(lr){2-3} \cmidrule(lr){4-5}
\multicolumn{1}{l}{\textbf{Estrato Geográfico}} & \textbf{Univariado} & \textbf{Multivariado} & \textbf{Univariado} & \textbf{Multivariado} \\
\midrule
01 - Belo Horizonte & 11,87 & 29,03 & -0,16 & 0,37 \\
02 - Entorno e Colar Metropolitano de BH & 6,85 & 26,42 & 0,19 & 1,29 \\
03 - Sul de Minas & 11,70 & 37,17 & -1,27 & -3,16 \\
04 - Triângulo Mineiro & 21,75 & 50,11 & -1,03 & -2,58 \\
05 - Zona da Mata & 16,86 & 20,48 & -1,41 & -1,02 \\
06 - Norte de Minas & 9,13 & 35,70 & -0,57 & -1,62 \\
07 - Vale do Rio Doce & 6,15 & 25,04 & -1,37 & -3,69 \\
08 - Central & 11,67 & 37,56 & -0,85 & -3,48 \\
\midrule
\textbf{Média} & \textbf{12,00} & \textbf{32,69} & & \\
\bottomrule
\end{tabular}%
}
}
\fonte{Elaboração própria, com base nos dados da PNAD Contínua. Nota: as oito primeiras observações de cada série foram descartadas do cálculo, em razão do período de estabilização do filtro de Kalman.}
\end{table}
